\pagebreak
\section{Журнал выполнения}
\begin{enumerate}
    \item Классы символов Юникода и понижение регистра не пришлось реализовывать самому: в glibc 2.39 локаль \texttt{C.UTF-8} встроена, и \texttt{iswalnum}/\texttt{towlower} работают для всего Юникода. Это проверено на символах разных письменностей (латиница с диакритикой, кириллица, греческий, иврит, арабский, иероглифы). Проверка же показала, что комбинируемые знаки (U+0301), тайские тоновые знаки и вирама деванагари буквами не считаются --- см. раздел о неудачных токенах.
    \item Первые замеры (ещё на первом корпусе ЛР1 из 30\,000 статей) показали, что больше половины времени уходит не на токенизацию, а на чтение файлов: даже \texttt{cat} всех файлов в \texttt{/dev/null} занимает около 3,7~с, из них 2,1~с --- системное время. Поэтому утилита считает время чтения и время токенизации раздельно, а скрипт замеров дополнительно измеряет тот же корпус, склеенный в один файл.
    \item Время одного и того же прогона колебалось в 1,5--2 раза. Процессор гибридный (Intel Core Ultra 7 255H: P-, E- и LP E-ядра с максимальной частотой 5,1, 4,4 и 2,5~ГГц), и планировщик ставит программу на любое из них; кроме того, часть замеров совпала с посторонней задачей, загрузившей все 16 ядер. Поэтому каждый размер измеряется несколько раз (размеры чередуются по кругу, между запусками пауза 0,5~с) и берётся минимальное время. На корпусе в 1 млн документов один круг всех размеров идёт около 5 минут, поэтому повторов 3 (на 30\,000 было 5); итоговые замеры сделаны при средней загрузке около 1,4.
    \item В частотном списке токенов нашлись склеенные слова вида \texttt{coordinates49}, \texttt{websitewww}, \texttt{passengersyear050}. Причина --- выделение текста в ЛР1: ячейки таблиц и инфобоксов (\texttt{<th>}, \texttt{<td>}) не отделялись друг от друга. Это ошибка выделения текста, а не правил токенизации. В корпусе на 1 млн документов таких склеек нет: статьи Wikipedia и Wikivoyage взяты из вики-разметки, где шаблоны инфобоксов вырезаются целиком.
    \item Корпус заменён на 1\,000\,000 документов четырёх источников (ЛР1), все замеры и частоты в отчёте пересчитаны на нём; токенизатор не менялся.
\end{enumerate}

\section{План тестирования}
Тесты --- юнит-тесты на C++ (\texttt{tests/test\_tokenizer.cpp}) на собственном минимальном каркасе \texttt{lab3/tests/unittest.h} (макросы \texttt{TEST} и \texttt{CHECK}, без сторонних библиотек). Они вызывают функции программы напрямую, а последний тест запускает утилиты командной строки как отдельные процессы. Сборка и запуск --- \texttt{make test}. Эталонные значения в тестах считаются независимо от проверяемого кода, для этого используются контейнеры STL (\texttt{std::map}, \texttt{std::set}) --- это код тестов, а не поисковой системы.

\begin{center}
\small
\begin{tabular}{|c|p{0.3\textwidth}|p{0.58\textwidth}|}
\hline
\textnumero & Тест & Что проверяется \\
\hline
1 & Декодирование UTF-8 & Последовательности из 1--4 байт дают правильный код символа и длину. \\
\hline
2 & Некорректный UTF-8 & Байт продолжения без начала, \texttt{0xFF}, оборванная, избыточная запись, суррогат, код больше U+10FFFF --- U+FFFD и длина 1. \\
\hline
3 & Кодирование UTF-8 & 12 кодов на границах длин 1--4 байта: декодирование закодированного символа возвращает тот же код. \\
\hline
4 & Английский текст & Знаки препинания и скобки --- разделители, регистр понижается, числа --- токены. \\
\hline
5 & Дефис и апостроф & \texttt{Hewlett-Packard}, \texttt{state-of-the-art}, \texttt{don't}, \texttt{Lviv’s} (апостроф U+2019) делятся на части. \\
\hline
6 & Числа & \texttt{3.14}, \texttt{1,000}, \texttt{B-52} делятся; \texttt{10am}, \texttt{19th} --- один токен. \\
\hline
7 & Юникод & Диакритика, кириллица, греческий остаются в токене, регистр понижается и у них (\texttt{ÉCOLE} $\to$ \texttt{école}, \texttt{İ} $\to$ \texttt{i}). \\
\hline
8 & Пробельные и типографские символы & Табуляция, перевод строки, неразрывный пробел, тире, кавычки-ёлочки, многоточие --- разделители. \\
\hline
9 & Нет токенов & Пустой текст и текст только из пробелов и знаков препинания. \\
\hline
10 & Иероглифы & Текст без пробелов становится одним токеном (известный недостаток, тест фиксирует поведение). \\
\hline
11 & Битые байты в тексте & Некорректный UTF-8 в середине и в конце текста --- разделитель, обработка не падает. \\
\hline
12 & Длина токена & Функция токенизации сообщает длину в символах, а не в байтах (\texttt{école} --- 5 символов, 6 байт). \\
\hline
13 & Статистика по файлам & Два файла с известным содержимым: число файлов, байтов, токенов и средняя длина совпадают с посчитанными вручную; несуществующий файл --- ошибка, статистика не меняется. \\
\hline
14 & Утилита \texttt{build/tokenizer} & Вывод токенов по одному на строку, режим \texttt{-s} по списку файлов, коды возврата при несуществующем файле (1) и неверном ключе (2). \\
\hline
\end{tabular}
\end{center}

\section{Результаты}
\subsection{Автотесты}
Пройдены все 14 тестов. Журнал прогона (\texttt{results/tests.log}; символы, которых нет в шрифте отчёта, записаны как \texttt{\textbackslash uXXXX}):

{\scriptsize
\VerbatimInput[breaklines, breakanywhere]{../results/tests.log}
}

\subsection{Количество и длина токенов}
Весь корпус (\texttt{results/stats.tsv}):
\begin{center}
\small
\begin{tabular}{|l|r|r|r|r|r|}
\hline
Источник & Док. & Текст, МиБ & Токенов & Токенов на док. & Ср. длина токена \\
\hline
Wikivoyage & 34\,818 & 260,6 & 46\,194\,627 & 1\,327 & 4,69 \\
\hline
Wikipedia & 620\,164 & 5\,011,8 & 865\,224\,622 & 1\,395 & 4,90 \\
\hline
Travel SE & 48\,768 & 97,3 & 18\,531\,043 & 380 & 4,37 \\
\hline
Reddit & 296\,250 & 333,4 & 65\,510\,878 & 221 & 4,16 \\
\hline
Итого & 1\,000\,000 & 5\,703,2 & 995\,461\,170 & 995 & 4,83 \\
\hline
\end{tabular}

\end{center}
Длина токена --- в символах. Самые длинные токены --- в статьях Wikipedia (4,90), короче --- в путеводителе Wikivoyage (4,69), самые короткие --- в постах Reddit (4,16) и вопросах Travel SE (4,37): разговорный язык проще энциклопедического. 4,4\,\% токенов --- односимвольные, в основном это осколки слов (см. ниже), 4,2\,\% --- числа.

\subsection{Время работы}
Время на первых $N$ документах корпуса (\texttt{results/bench.tsv}; минимум из 3 запусков). \enquote{Всего} --- время программы от начала до конца, в т.\,ч. чтение файлов; последние два столбца --- время на килобайт входного текста, всего и только на токенизацию. Последняя строка --- тот же корпус, склеенный в один файл.
\begin{center}
\footnotesize
\begin{tabular}{|l|r|r|r|r|r|r|}
\hline
Файлов & МиБ & Всего, с & Чтение, с & Токениз., с & мкс/КиБ & мкс/КиБ (токениз.) \\
\hline
1\,000 & 5,7 & 0,12 & 0,08 & 0,03 & 20,4 & 5,3 \\
\hline
3\,000 & 15,3 & 0,27 & 0,19 & 0,08 & 17,2 & 5,0 \\
\hline
10\,000 & 46,2 & 1,01 & 0,74 & 0,26 & 21,4 & 5,4 \\
\hline
30\,000 & 137,0 & 2,43 & 1,73 & 0,67 & 17,3 & 4,8 \\
\hline
100\,000 & 455,5 & 7,62 & 5,27 & 2,24 & 16,3 & 4,8 \\
\hline
300\,000 & 1\,372,4 & 19,61 & 13,11 & 6,20 & 14,0 & 4,4 \\
\hline
1\,000\,000 & 5\,703,2 & 164,97 & 129,22 & 34,15 & 28,2 & 5,8 \\
\hline
1 (склеенный) & 5\,703,2 & 36,97 & 10,11 & 26,87 & 6,3 & 4,6 \\
\hline
\end{tabular}

\end{center}

Собственно токенизация растёт линейно с объёмом: на всех размерах 4,4--5,8~мкс/КиБ, без тренда. Так и должно быть: каждый символ обрабатывается за постоянное время, состояние между документами не накапливается. Общее удельное время до 300\,000 файлов тоже постоянно (14--21~мкс/КиБ, колебания --- шум замеров, см. журнал), а на всём корпусе вырастает до 28~мкс/КиБ из-за чтения: миллион файлов читается 129~с, 130~мкс на файл против 44~мкс на 300\,000 файлов. Видимо, весь корпус (5,6~ГиБ мелкими файлами) вместе со склеенной копией для последней строки уже не держится в кэше ОС, и часть файлов читается с диска. Весь корпус обрабатывается за 165~с (35~МиБ/с), из них собственно токенизация --- 34~с (около 5,8~мкс/КиБ, 170~МиБ/с).

\subsection{Оптимальна ли скорость и как её ускорить}
Скорость не оптимальна: предел --- скорость чтения данных из памяти (единицы ГиБ/с), а программа работает в десятки раз медленнее.
\begin{enumerate}
    \item Около 80\,\% времени --- открытие и чтение миллиона мелких файлов (системные вызовы и диск). Тот же корпус одним файлом читается за 10~с вместо 129~с, и весь прогон занимает 37~с вместо 165~с (последняя строка таблицы; время токенизации при этом почти то же --- код тот же). Ускорение: хранить корпус в нескольких больших файлах, читать через \texttt{mmap}.
    \item Документы независимы, поэтому их можно токенизировать параллельно в нескольких потоках; на 16-ядерной машине это ускорит работу в разы, пока не упрётся в память или диск.
    \item Сама токенизация --- около 5~нс на байт. Её можно ускорить: проверять класс символа по заранее построенной таблице (битовая карта на 65\,536 символов) вместо вызова \texttt{iswalnum}; не копировать токены из ASCII-символов в \texttt{std::string}, а передавать указатель на участок буфера с понижением регистра на месте; искать границы токенов в ASCII-тексте SIMD-инструкциями по 16--32 байта.
\end{enumerate}
Собственно токенизация миллиона документов (полминуты) --- небольшая часть построения индекса (ЛР7), основное время --- чтение файлов, поэтому эти ускорения в работе не делались; первым стоило бы склеить корпус в большие файлы.

\subsection{Неудачно выделенные токены}
Частотный список токенов корпуса построен тем же токенизатором (скрипт ЛР5): 3\,974\,421 разный токен. В таблице --- найденные в нём неудачные случаи, число в скобках --- сколько раз токен встретился в корпусе.

{\small
\begin{longtable}{|>{\raggedright\arraybackslash}p{0.2\textwidth}|>{\raggedright\arraybackslash}p{0.19\textwidth}|>{\raggedright\arraybackslash}p{0.2\textwidth}|>{\raggedright\arraybackslash}p{0.27\textwidth}|}
\hline
Текст & Токены & Проблема & Как исправить \\
\hline
\endhead
\texttt{Lviv's}, \texttt{Lviv’s} & \texttt{lviv}, \texttt{s} & Притяжательное \texttt{'s} даёт токен \texttt{s} (6\,753\,822, 13-й по частоте) & Апостроф между буквами оставлять внутри токена, окончание \texttt{'s} отбрасывать при нормализации \\
\hline
\texttt{don't} & \texttt{don}, \texttt{t} & Слово разрушено, \texttt{don} совпадает с именем Don, \texttt{t} (698\,944) & Апостроф между буквами --- часть токена: \texttt{don't} \\
\hline
\texttt{Hewlett-Packard}, \texttt{e-mail} & \texttt{hewlett}, \texttt{packard}; \texttt{e}, \texttt{mail} & Составное слово разбито; \texttt{e-mail} и \texttt{email} (17\,470) --- разные токены & Индексировать и части, и слово целиком без дефиса (\texttt{email}) \\
\hline
\texttt{U.S.}, \texttt{AT\&T}, \texttt{C++} & \texttt{u}, \texttt{s}; \texttt{at}, \texttt{t}; \texttt{c} & Сокращения и названия разрушены & Шаблон \enquote{буква, точка, буква, точка} $\to$ \texttt{us}; словарь исключений (\texttt{AT\&T}, \texttt{C++}) \\
\hline
\texttt{3,000}, \texttt{49.84167} & \texttt{3}, \texttt{000}; \texttt{49}, \texttt{84167} & Число разбито: токен \texttt{000} (693\,257) & Цифры, разделённые одной точкой или запятой, --- один токен; разделитель тысяч убирать (\texttt{3000}) \\
\hline
\texttt{49°50$'$30$''$N} (координаты) & \texttt{49}, \texttt{50}, \texttt{30}, \texttt{n} & Мусорные числа и буква \texttt{n} (81\,924) & Распознавать координаты шаблоном и не индексировать или хранить отдельно \\
\hline
\texttt{10am} и \texttt{10 am} & \texttt{10am}; \texttt{10}, \texttt{am} & Одно и то же --- разные токены (\texttt{10am}: 19\,175) & Делить на границе цифр и букв, одинаково в документах и запросах \\
\hline
\texttt{www.\allowbreak voodoofestival.\allowbreak com} & \texttt{www}, \texttt{voodoofestival}, \texttt{com} & Адрес сайта разбит, \texttt{com} (53\,489) & Шаблоны URL и e-mail --- один токен (и, возможно, его части) \\
\hline
\texttt{café} и \texttt{cafe}, \texttt{São} и \texttt{Sao}, \texttt{Straße} и \texttt{Strasse} & разные токены & \texttt{café} 22\,061 против \texttt{cafe} 30\,907: запрос без диакритики не найдёт 42\,\% & Нормализация: снятие диакритики, \texttt{ß} $\to$ \texttt{ss}, совместимые символы (NFKC: лигатура \texttt{fi} U+FB01 $\to$ \texttt{fi}) \\
\hline
Тайское слово \enquote{songsoem} (тоновый знак U+0E48 после первой буквы), хинди \enquote{hindi} (вирама U+094D), IPA с ударением U+0301 внутри слова & 2 токена: часть слова до знака и после & Комбинируемые знаки не считаются буквами, слово рвётся & Знаки категорий Юникода Mn и Mc после буквы считать частью токена \\
\hline
\texttt{SpaceShipOne} с мягкими переносами U+00AD & \texttt{space}, \texttt{ship}, \texttt{one} & Невидимый символ рвёт слово & Удалять U+00AD, U+200C, U+200D до токенизации \\
\hline
Тайские и китайские фразы без пробелов & один токен до 75 символов & Фраза целиком --- один токен, по слову не найти (162\,391 вхождение, 0,02\,\%) & Сегментация по словарю или биграммы символов \\
\hline
\texttt{Coordinates49°50$'$}, \texttt{Websitewww.\ldots}, \texttt{Elevation AMSL13 ft} & \texttt{coordinates49}, \texttt{websitewww}, \texttt{amsl13} & Склеены соседние ячейки инфобокса (в корпусе 30 тыс. из HTML; в миллионе инфобоксы вырезаются из вики-разметки, склеек нет) & Исправить выделение текста в ЛР1: \texttt{<td>}, \texttt{<th>} --- разделители \\
\hline
\end{longtable}
}

\subsection{Достоинства и недостатки метода}
Достоинства:
\begin{itemize}
    \item одно простое правило, один проход по тексту, без словарей и регулярных выражений; время линейно от объёма;
    \item работает для любой письменности с пробелами между словами (латиница, кириллица, греческий и т.\,д.), а не только для английского;
    \item одинаково применяется к документам и запросам: запрос \enquote{Hewlett-Packard} разбивается так же, как текст, и документ находится;
    \item некорректный UTF-8 не ломает обработку.
\end{itemize}
Недостатки:
\begin{itemize}
    \item разрушаются единицы со знаками внутри: сокращения, числа с разделителями, URL, координаты, слова с апострофом и дефисом; появляются мусорные короткие токены;
    \item нет нормализации диакритики и совместимых символов;
    \item не подходит для языков без пробелов и рвёт слова с комбинируемыми знаками.
\end{itemize}

\section{Выводы}
Написан токенизатор на C++: токен --- непрерывная последовательность букв и цифр Юникода, регистр понижается. В корпусе из 1\,000\,000 документов 995~млн токенов средней длиной 4,83 символа. Время токенизации линейно зависит от объёма (около 5~мкс/КиБ); весь корпус (5,6~ГиБ) обрабатывается за 165~с, из которых токенизация --- 34~с, остальное --- чтение миллиона файлов.

Правило простое и быстрое, но рвёт на части сокращения, числа, URL, слова с апострофом, дефисом и комбинируемыми знаками и не объединяет написания с диакритикой и без. Исправления: отдельные шаблоны для чисел, URL, сокращений; апостроф и дефис между буквами внутри токена; комбинируемые знаки --- часть слова; нормализация NFKC и снятие диакритики. Скорость не оптимальна: её ограничивает чтение мелких файлов, ускорить можно склейкой корпуса в большие файлы, параллельной обработкой и табличной проверкой классов символов.
